\chapter*{Der Halbgruppenbeweis}
\hypertarget{power.reading}{}
\addcontentsline{toc}{chapter}{Warum die Potenzhalbgruppe eine Halbgruppe ist}

Sei \((A,\star)\) eine Halbgruppe. Ihre nichtleeren Teilmengen bilden
die Menge
\[
 \PowNE{A}=\{X\subseteq A\mid X\ne\varnothing\}.
\]
Für \(X,Y\in\PowNE{A}\) ist das \emph{Mengenprodukt} nach
\FormulaRefAuto{PowerSemigroupProductDef}[def] definiert durch
\[
 X\star_{\mathcal P}Y
   \coloneqq\{x\star y\mid x\in X,\ y\in Y\}.
\]
Die rechte Seite ist die Termbildschreibweise mit zwei Indizes aus
\FormulaRefAuto{B27TermImageDef}[def].
Mit diesem Mengenprodukt ist \(\PowNE{A}\) eine Halbgruppe, die
\emph{Potenzhalbgruppe} von \((A,\star)\)
(\FormulaRefAuto{NonemptyPowerSemigroup}[thm]).

Für nichtleere Teilmengen sichert
\FormulaRefAuto{BinaryOperationImageNonemptyPowerset}[thm] unmittelbar die Abgeschlossenheit:
\[
 \begin{aligned}
 X\star_{\mathcal P}Y
 &\overset{\text{\FormulaRefAuto{PowerSemigroupProductDef}[def]}}{=}
   \{x\star y\mid x\in X,\ y\in Y\}\\[5pt]
 &\overset{\text{\FormulaRefAuto{BinaryOperationImageNonemptyPowerset}[thm]}}{\in}
   \PowNE{A}.
 \end{aligned}
\]
Der zitierte Satz verbindet die Zugehörigkeit zur Potenzmenge aus
\FormulaRefAuto{BinaryOperationImagePowerset}[thm] mit der Nichtleerheit aus
\FormulaRefAuto{BinaryTermImageNonempty}[thm].

Für \(X,Y,Z\in\PowNE{A}\) folgt die Assoziativität aus der Kette
\[
 \begin{aligned}
 &X\star_{\mathcal P}(Y\star_{\mathcal P}Z)\\[3pt]
 &\quad\overset{\text{\FormulaRefAuto{TermImageCompositionRightFromEquality}[thm]}}{=}
   \{x\star(y\star z)\mid x\in X,\ y\in Y,\ z\in Z\}\\[5pt]
 &\quad\overset{\substack{\text{\FormulaRefAuto{SemigroupAssociativityAxiom}[ax]}\\
                         \text{\FormulaRefAuto{TripleTermImageEqualityIntroductionRule}[thm]}}}{=}
   \{(x\star y)\star z\mid x\in X,\ y\in Y,\ z\in Z\}\\[5pt]
 &\quad\overset{\text{\FormulaRefAuto{TermImageCompositionLeftFromEquality}[thm]}}{=}
   (X\star_{\mathcal P}Y)\star_{\mathcal P}Z.
 \end{aligned}
\]
Zusammen mit der Abgeschlossenheit ist nach
\FormulaRefAuto{SemigroupDef}[def] die Halbgruppeneigenschaft bewiesen.

\clearpage
\section*{Ein Beispiel mit drei Elementen}

\begin{PowerReadingExample}{Addition modulo drei und ihre Potenzhalbgruppe}
Auf \(A=\{0,1,2\}\) sei \(x\star y\) der Rest der Summe
\(x+y\) bei Division durch drei. Diese Operation ist abgeschlossen
und assoziativ: Beide Klammerungen ergeben den Rest derselben Summe
von drei Zahlen.

Die Potenzhalbgruppe hat sieben Elemente: die drei Einermengen, die
drei zweielementigen Teilmengen und \(A\) selbst. Ihre vollständige
Multiplikationstabelle lautet:
\[
 {\setlength{\arraycolsep}{3.5pt}
 \begin{array}{c|ccccccc}
 \star_{\mathcal P}&\{0\}&\{1\}&\{2\}&\{0,1\}&\{0,2\}&\{1,2\}&A\\\hline
 \{0\}&\{0\}&\{1\}&\{2\}&\{0,1\}&\{0,2\}&\{1,2\}&A\\
 \{1\}&\{1\}&\{2\}&\{0\}&\{1,2\}&\{0,1\}&\{0,2\}&A\\
 \{2\}&\{2\}&\{0\}&\{1\}&\{0,2\}&\{1,2\}&\{0,1\}&A\\
 \{0,1\}&\{0,1\}&\{1,2\}&\{0,2\}&A&A&A&A\\
 \{0,2\}&\{0,2\}&\{0,1\}&\{1,2\}&A&A&A&A\\
 \{1,2\}&\{1,2\}&\{0,2\}&\{0,1\}&A&A&A&A\\
 A&A&A&A&A&A&A&A
 \end{array}}
\]
Der Zeilenfaktor steht links, der Spaltenfaktor rechts.

Für \(X=\{0,1\}\), \(Y=\{1,2\}\) entstehen beispielsweise
die Produktwerte
\[
 \begin{array}{c|cc}
 \star&1&2\\\hline
 0&1&2\\
 1&2&0
 \end{array}
 \qquad\text{also}\qquad X\star_{\mathcal P}Y=A.
\]
Der Wert \(2\) entsteht zweimal und gehört trotzdem nur einmal
zur Ergebnis\emph{menge}.

Mit \(Z=\{1\}\) ist \(Y\star_{\mathcal P}Z=\{0,2\}\).
Die beiden Klammerungen liefern
\[
 \begin{aligned}
 X\star_{\mathcal P}(Y\star_{\mathcal P}Z)
   &=\{0,1\}\star_{\mathcal P}\{0,2\}=A,\\
 (X\star_{\mathcal P}Y)\star_{\mathcal P}Z
   &=A\star_{\mathcal P}\{1\}=A.
 \end{aligned}
\]
Die Zwischenmengen sind verschieden, die Endergebnisse gleich.
\end{PowerReadingExample}

\clearpage
\section*{Historischer Abriss}

Die Mengenmultiplikation ist älter als die Bezeichnung
\emph{Potenzhalbgruppe}. Paul Dubreil beschreibt in seiner Arbeit
\emph{Contribution à la théorie des demi-groupes. III} von 1953
bereits ausdrücklich die Halbgruppe der Teilmengen eines Halbgruppenträgers
(Abschnitt~1, S.~290). Sein Produkt besteht aus allen Produkten
\(ab\) mit \(a\) aus der ersten und \(b\) aus der zweiten Menge.
Dubreil nimmt zusätzlich die leere Menge auf; sie ist ein Nullelement.
Die hier betrachteten nichtleeren Teilmengen bilden darin eine
Unterhalbgruppe.
\href{https://doi.org/10.24033/bsmf.1444}{Originalarbeit: Bull. Soc. Math. France 81 (1953), 289--306.}

1967 veröffentlichten Takayuki Tamura und John Shafer die Arbeit
\emph{Power semigroups}, Math. Japon. 12, S.~25--32, einen frühen
Beitrag zur systematischen Untersuchung dieser Konstruktion.
Tamuras Bericht \emph{Recent Study of Semigroups},
RIMS K\^oky\^uroku 31 (1967), S.~1--35, enthält in Problem~10
auf S.~35 genau die hier betrachtete Konstruktion auf den nichtleeren
Teilmengen und hält ausdrücklich fest, dass sie eine Halbgruppe ergibt.
\href{https://repository.kulib.kyoto-u.ac.jp/server/api/core/bitstreams/d91e4d3f-686c-448c-b6a8-318edf723d0b/content\#page=35}{Originalscan: Problem~10, S.~35.}
Tamura schließt daran die Frage an, ob isomorphe Potenzhalbgruppen
auch isomorphe Grundhalbgruppen erzwingen. Der elementare Nachweis
der Halbgruppeneigenschaft bildet damit den Ausgangspunkt für die
späteren Struktur- und Rekonstruktionsfragen.

\medskip\noindent\PowerMainLink{}\enspace\textperiodcentered\enspace\PowerProofsLink{}
